% GENERATED FILE. Edit canonical lessons, then run npm run generate:books.
% canonical-source-hash: fnv1a64:68b88c4fdfb1a249
% canonical-lessons: JA-C189-hakaru, JA-C189-todokeru, JA-C189-uketoru, JA-C189-watasu, JA-C189-omiyage, JA-R189-review-184-189

\chapter{Weighing and Handing Over}
\label{ch:ja-weighing-and-handing-over}
\hlchaptermodality{189}

\begin{chapteropening}
\textbf{By the end of this chapter:} \emph{I can weigh, deliver and hand over parcels.}

Complete the last lesson of chapter 189: I can weigh, deliver and hand over parcels.
\end{chapteropening}

\section[\texorpdfstring{\ja{はかる} (hakaru)}{はかる (hakaru)}]{\ja{はかる} (hakaru) — to weigh; to measure}

\label{lesson:JA-C189-hakaru}



\begin{quote}
Before the new one: say the Japanese for to send, then the Japanese for registered mail.
\end{quote}

\subsection*{You'll want to know: \ja{はかる}}
\textbf{\ja{はかる}} — \emph{hakaru} — \textquotedblleft{}to weigh; to measure\textquotedblright{}.

\begin{quote}
\textbf{\ja{ゆうびんきょくで} \ja{にもつを} \ja{はかります。}} — \emph{yūbinkyoku de nimotsu o hakarimasu} — \textquotedblleft{}They weigh the package at the post office.\textquotedblright{}
\end{quote}

Polite \textbf{\ja{はかります}}, \emph{hakarimasu}; a godan verb.

\begin{cousinweb}
\textbf{\ja{はかる}} (\emph{hakaru}) is a native Japanese word (wago), inherited rather than borrowed from Chinese; its kanji was attached later as a kun'yomi (JMdict, KANJIDIC2).
\end{cousinweb}

\begin{grammarlens}[title={What you've built}]
One more word for everyday life at A2.
\end{grammarlens}

\subsection*{Your turn}
\begin{itemize}
\raggedright
  \item \emph{Say it:} \emph{hakaru}
  \item \emph{Say it:} \emph{hakaru}, once more
  \item \emph{Say it:} \emph{yūbinkyoku de nimotsu o hakarimasu} — They weigh the package at the post office.
  \item \emph{Recall:} say the Japanese for to send, then the Japanese for registered mail, then say \emph{hakaru} again
\end{itemize}

\begin{tcolorbox}[breakable,colback=teal!4,colframe=teal!35!black,title={Before you move on}]
What does \ja{はかる} mean? (\textquotedblleft{}To weigh; to measure\textquotedblright{}.) Say it once more.
\end{tcolorbox}
\section[\texorpdfstring{\ja{とどける} (todokeru)}{とどける (todokeru)}]{\ja{とどける} (todokeru) — to deliver; to report}

\label{lesson:JA-C189-todokeru}



\begin{quote}
Before the new one: say the Japanese for registered mail, then the Japanese for to weigh.
\end{quote}

\subsection*{You'll want to know: \ja{とどける}}
\textbf{\ja{とどける}} — \emph{todokeru} — \textquotedblleft{}to deliver; to report\textquotedblright{}.

\begin{quote}
\textbf{\ja{あした} \ja{にもつを} \ja{とどけます。}} — \emph{ashita nimotsu o todokemasu} — \textquotedblleft{}I will deliver the package tomorrow.\textquotedblright{}
\end{quote}

Polite \textbf{\ja{とどけます}}, \emph{todokemasu}; an ichidan verb. Its intransitive partner \textbf{\ja{とどく}} is taught already.

\begin{cousinweb}
\textbf{\ja{とどける}} (\emph{todokeru}) is a native Japanese word (wago), inherited rather than borrowed from Chinese; its kanji was attached later as a kun'yomi, and it shares its stem with \emph{todoku}, \textquotedblleft{}to reach, to arrive\textquotedblright{} (JMdict, KANJIDIC2).
\end{cousinweb}

\begin{grammarlens}[title={What you've built}]
One more word for everyday life at A2.
\end{grammarlens}

\subsection*{Your turn}
\begin{itemize}
\raggedright
  \item \emph{Say it:} \emph{todokeru}
  \item \emph{Say it:} \emph{todokeru}, once more
  \item \emph{Say it:} \emph{ashita nimotsu o todokemasu} — I will deliver the package tomorrow.
  \item \emph{Recall:} say the Japanese for registered mail, then the Japanese for to weigh, then say \emph{todokeru} again
\end{itemize}

\begin{tcolorbox}[breakable,colback=teal!4,colframe=teal!35!black,title={Before you move on}]
What does \ja{とどける} mean? (\textquotedblleft{}To deliver; to report\textquotedblright{}.) Say it once more.
\end{tcolorbox}
\section[\texorpdfstring{\ja{うけとる} (uketoru)}{うけとる (uketoru)}]{\ja{うけとる} (uketoru) — to accept; to receive, to take delivery of}

\label{lesson:JA-C189-uketoru}



\begin{quote}
Before the new one: say the Japanese for to weigh, then the Japanese for to deliver.
\end{quote}

\subsection*{You'll want to know: \ja{うけとる}}
\textbf{\ja{うけとる}} — \emph{uketoru} — \textquotedblleft{}to accept; to receive, to take delivery of\textquotedblright{}.

\begin{quote}
\textbf{\ja{うけつけで} \ja{にもつを} \ja{うけとります。}} — \emph{uketsuke de nimotsu o uketorimasu} — \textquotedblleft{}I collect the package at reception.\textquotedblright{}
\end{quote}

Polite \textbf{\ja{うけとります}}, \emph{uketorimasu}; a godan verb.

\begin{cousinweb}
\textbf{\ja{うけとる}} (\emph{uketoru}) is native Japanese, built from \emph{uke}, the stem of \emph{ukeru}, \textquotedblleft{}to receive\textquotedblright{}, and \emph{toru}, \textquotedblleft{}to take\textquotedblright{}, both inherited rather than borrowed (JMdict, KANJIDIC2).
\end{cousinweb}

\begin{grammarlens}[title={What you've built}]
One more word for everyday life at A2.
\end{grammarlens}

\subsection*{Your turn}
\begin{itemize}
\raggedright
  \item \emph{Say it:} \emph{uketoru}
  \item \emph{Say it:} \emph{uketoru}, once more
  \item \emph{Say it:} \emph{uketsuke de nimotsu o uketorimasu} — I collect the package at reception.
  \item \emph{Recall:} say the Japanese for to weigh, then the Japanese for to deliver, then say \emph{uketoru} again
\end{itemize}

\begin{tcolorbox}[breakable,colback=teal!4,colframe=teal!35!black,title={Before you move on}]
What does \ja{うけとる} mean? (\textquotedblleft{}To accept; to receive, to take delivery of\textquotedblright{}.) Say it once more.
\end{tcolorbox}
\section[\texorpdfstring{\ja{わたす} (watasu)}{わたす (watasu)}]{\ja{わたす} (watasu) — to hand over; to pass, to give}

\label{lesson:JA-C189-watasu}



\begin{quote}
Before the new one: say the Japanese for to deliver, then the Japanese for to accept.
\end{quote}

\subsection*{You'll want to know: \ja{わたす}}
\textbf{\ja{わたす}} — \emph{watasu} — \textquotedblleft{}to hand over; to pass, to give\textquotedblright{}.

\begin{quote}
\textbf{\ja{てんいんに} \ja{おかねを} \ja{わたします。}} — \emph{ten'in ni okane o watashimasu} — \textquotedblleft{}I hand the money to the shop assistant.\textquotedblright{}
\end{quote}

Polite \textbf{\ja{わたします}}, \emph{watashimasu}; a godan verb.

\begin{cousinweb}
\textbf{\ja{わたす}} (\emph{watasu}) is a native Japanese word (wago), inherited rather than borrowed from Chinese; its kanji was attached later as a kun'yomi (JMdict, KANJIDIC2).
\end{cousinweb}

\begin{grammarlens}[title={What you've built}]
One more word for everyday life at A2.
\end{grammarlens}

\subsection*{Your turn}
\begin{itemize}
\raggedright
  \item \emph{Say it:} \emph{watasu}
  \item \emph{Say it:} \emph{watasu}, once more
  \item \emph{Say it:} \emph{ten'in ni okane o watashimasu} — I hand the money to the shop assistant.
  \item \emph{Recall:} say the Japanese for to deliver, then the Japanese for to accept, then say \emph{watasu} again
\end{itemize}

\begin{tcolorbox}[breakable,colback=teal!4,colframe=teal!35!black,title={Before you move on}]
What does \ja{わたす} mean? (\textquotedblleft{}To hand over; to pass, to give\textquotedblright{}.) Say it once more.
\end{tcolorbox}
\section[\texorpdfstring{\ja{おみやげ} (omiyage)}{おみやげ (omiyage)}]{\ja{おみやげ} (omiyage) — souvenir; gift brought back from a trip}

\label{lesson:JA-C189-omiyage}



\begin{quote}
Before the new one: say the Japanese for to accept, then the Japanese for to hand over.
\end{quote}

\subsection*{You'll want to know: \ja{おみやげ}}
\textbf{\ja{おみやげ}} — \emph{omiyage} — \textquotedblleft{}souvenir; gift brought back from a trip\textquotedblright{}.

\begin{quote}
\textbf{\ja{ともだちに} \ja{おみやげを} \ja{かいます。}} — \emph{tomodachi ni omiyage o kaimasu} — \textquotedblleft{}I buy a souvenir for a friend.\textquotedblright{}
\end{quote}

\begin{cousinweb}
\textbf{\ja{おみやげ}} (\emph{omiyage}) is built from polite \emph{o-} and native \emph{miyage}, written with two kanji for \textquotedblleft{}local product\textquotedblright{} that are read as a whole; the older origin of \emph{miyage} is not settled (JMdict, KANJIDIC2).
\end{cousinweb}

\begin{grammarlens}[title={What you've built}]
One more word for everyday life at A2.
\end{grammarlens}

\subsection*{Your turn}
\begin{itemize}
\raggedright
  \item \emph{Say it:} \emph{omiyage}
  \item \emph{Say it:} \emph{omiyage}, once more
  \item \emph{Say it:} \emph{tomodachi ni omiyage o kaimasu} — I buy a souvenir for a friend.
  \item \emph{Recall:} say the Japanese for to accept, then the Japanese for to hand over, then say \emph{omiyage} again
\end{itemize}

\begin{tcolorbox}[breakable,colback=teal!4,colframe=teal!35!black,title={Before you move on}]
What does \ja{おみやげ} mean? (\textquotedblleft{}Souvenir; gift brought back from a trip\textquotedblright{}.) Say it once more.
\end{tcolorbox}
\section[(dialogue)]{Words from chapters 184–189 — a review}

\label{lesson:JA-R189-review-184-189}



\begin{quote}
Say the Japanese for goods and queue.
\end{quote}

\subsection*{Your turn}
Say these aloud:

\begin{itemize}
\raggedright
  \item goods, queue, to exchange, all of it, exactly
  \item discount, bargain sale, sold out, to cost, to be enough
  \item to be in surplus, savings, to save up, fee, tax
  \item to deposit, to withdraw, bankbook, to borrow, money exchange
  \item post office, stamp, postcard, to send, registered mail
  \item to weigh, to deliver, to accept, to hand over, souvenir
\end{itemize}

\begin{tcolorbox}[breakable,colback=teal!4,colframe=teal!35!black,title={Before you move on}]
Goods? (\textbf{\ja{しなもの}}, \emph{shinamono}.) Queue? (\textbf{\ja{れつ}}, \emph{retsu}.) To exchange? (\textbf{\ja{とりかえる}}, \emph{torikaeru}.) All of it? (\textbf{\ja{ぜんぶ}}, \emph{zenbu}.) Exactly? (\textbf{\ja{ちょうど}}, \emph{chōdo}.) Discount? (\textbf{\ja{わりびき}}, \emph{waribiki}.) Bargain sale? (\textbf{\ja{やすうり}}, \emph{yasuuri}.) Sold out? (\textbf{\ja{うりきれ}}, \emph{urikire}.) To cost? (\textbf{\ja{かかる}}, \emph{kakaru}.) To be enough? (\textbf{\ja{たりる}}, \emph{tariru}.) To be in surplus? (\textbf{\ja{あまる}}, \emph{amaru}.) Savings? (\textbf{\ja{ちょきん}}, \emph{chokin}.) To save up? (\textbf{\ja{ためる}}, \emph{tameru}.) Fee? (\textbf{\ja{りょうきん}}, \emph{ryōkin}.) Tax? (\textbf{\ja{ぜいきん}}, \emph{zeikin}.) To deposit? (\textbf{\ja{あずける}}, \emph{azukeru}.) To withdraw? (\textbf{\ja{おろす}}, \emph{orosu}.) Bankbook? (\textbf{\ja{つうちょう}}, \emph{tsūchō}.) To borrow? (\textbf{\ja{かりる}}, \emph{kariru}.) Money exchange? (\textbf{\ja{りょうがえ}}, \emph{ryōgae}.) Post office? (\textbf{\ja{ゆうびんきょく}}, \emph{yūbinkyoku}.) Stamp? (\textbf{\ja{きって}}, \emph{kitte}.) Postcard? (\textbf{\ja{はがき}}, \emph{hagaki}.) To send? (\textbf{\ja{おくる}}, \emph{okuru}.) Registered mail? (\textbf{\ja{かきとめ}}, \emph{kakitome}.) To weigh? (\textbf{\ja{はかる}}, \emph{hakaru}.) To deliver? (\textbf{\ja{とどける}}, \emph{todokeru}.) To accept? (\textbf{\ja{うけとる}}, \emph{uketoru}.) To hand over? (\textbf{\ja{わたす}}, \emph{watasu}.) Souvenir? (\textbf{\ja{おみやげ}}, \emph{omiyage}.)
\end{tcolorbox}
